\chapter{Marches aléatoires et martingales}\label{ch:iv1:random-walks-and-martingales}

\begin{ivtrains}
Probabilités d'atteinte, durées et extrêmes des marches, généralement sans résoudre de récurrence~: repérer ce qui reste équitable, arrêter le processus et lire la réponse. Les questions suivantes demandent d'inventer soi-même la martingale ou d'expliquer pourquoi l'arrêt optionnel échoue.
\end{ivtrains}

\begin{question}[$\star$ \iqroles{QR, QT, Trader, Risque} \iqfirm{toutes} \iqlevel{junior}]\label{iq:iv1:random-walks-and-martingales:1}
Vous vous trouvez à la position 3 sur une droite numérotée de 0 à 10. Chaque pas vous déplace d'une case vers la gauche ou vers la droite avec la même probabilité, et vous vous arrêtez dès que vous atteignez 0 ou 10. Quelle est la probabilité que vous finissiez en 10, et combien de pas faites-vous en moyenne~?
\end{question}

\begin{question}[$\star$ \iqroles{QR, QT, Trader, Risque} \iqfirm{toutes} \iqlevel{junior}]\label{iq:iv1:random-walks-and-martingales:2}
Vous misez 1~\$ à la fois, à parité, et vous gagnez chaque pari avec la probabilité $0.6$. Vous partez avec 5~\$ et vous arrêtez dès que vous atteignez 10~\$ ou 0~\$. Quelle est la probabilité que vous atteigniez 10~\$, et combien de paris faut-il en moyenne~?
\end{question}

\begin{question}[$\star$ \iqroles{QT, Trader, Analyste, Risque} \iqfirm{toutes} \iqlevel{junior}]\label{iq:iv1:random-walks-and-martingales:3}
Vous jouez à un jeu équitable, en misant 1~\$ à chaque fois à parité, et vous vous arrêtez dès que vous êtes en gain de 1~\$ ou en perte de 9~\$. À quelle fréquence gagnez-vous le 1~\$~? Quel est votre gain moyen, et combien de paris dure une séance en moyenne~?
\end{question}

\begin{question}[$\star$ \iqroles{QR, QT, Trader} \iqfirm{toutes} \iqlevel{junior}]\label{iq:iv1:random-walks-and-martingales:4}
Une marche aléatoire part de 0 et se déplace de $+1$ ou $-1$ avec la même probabilité à chaque pas. Quelle est la probabilité qu'elle soit revenue en 0 après 10 pas~? À quelle distance de 0 se trouve-t-elle en moyenne après 100 pas~? Donnez une formule exacte et l'approximation rapide.
\end{question}

\begin{question}[$\star\star$ \iqroles{QR, QT, Trader, Risque} \iqfirm{toutes} \iqlevel{junior}]\label{iq:iv1:random-walks-and-martingales:5}
Une marche se déplace de $+1$ avec la probabilité $0.6$ et de $-1$ avec la probabilité $0.4$ à chaque pas. En moyenne, combien de temps faut-il pour aller de 0 à $+10$ pour la première fois~? Quelle identité vous donne ce résultat sans aucune équation~?
\end{question}

\begin{question}[$\star\star$ \iqroles{QR, QT, Risque} \iqfirm{toutes} \iqlevel{junior}]\label{iq:iv1:random-walks-and-martingales:6}
Une marche se déplace de $+1$ avec la probabilité $0.4$ et de $-1$ avec la probabilité $0.6$ à chaque pas, indéfiniment. Quelle est la probabilité qu'elle atteigne un jour $+10$~? Quelle est la valeur moyenne du point le plus haut qu'elle atteint jamais~?
\end{question}

\begin{question}[$\star\star$ \iqroles{QT, Trader, Risque, Analyste} \iqfirm{toutes} \iqlevel{junior}]\label{iq:iv1:random-walks-and-martingales:7}
Vous disposez de 1\,023~\$ et jouez à un jeu équitable à parité en doublant~: misez 1~\$, puis, après chaque perte, doublez la mise, en vous arrêtant dès que vous gagnez ou que vous ne pouvez plus miser le pari suivant. Quelle est la probabilité de terminer avec 1~\$ de gain~? Quel est votre gain espéré~? Si une victoire met fin à la séance, combien de paris faites-vous en moyenne~?
\end{question}

\begin{question}[$\star\star$ \iqroles{QR, QT, Risque} \iqfirm{hedge fund, teneur de marché} \iqlevel{junior}]\label{iq:iv1:random-walks-and-martingales:8}
Une marche équitable part de 0 et s'arrête la première fois qu'elle atteint $+1$. Elle atteint $+1$ presque sûrement, si bien que sa valeur finale vaut toujours 1, pourtant c'est un jeu équitable qui part de 0. Pourquoi cela ne contredit-il pas l'affirmation selon laquelle un jeu équitable ne peut pas être battu par une règle d'arrêt~? Quel est le temps d'arrêt médian, et comment se comporte la moyenne~?
\end{question}

\begin{question}[$\star\star$ \iqroles{QR, QT, QD} \iqfirm{toutes} \iqlevel{junior}]\label{iq:iv1:random-walks-and-martingales:9}
Une marche sur les entiers de 0 à 10 se déplace vers la gauche ou vers la droite avec la même probabilité, sauf qu'en 0 elle passe toujours en 1. En moyenne, combien de pas faut-il pour aller de 0 à 10 pour la première fois~? Et de 4 à 10~? Trouvez un processus qui reste équitable et qui donne directement ce résultat.
\end{question}

\begin{question}[$\star\star$ \iqroles{QR, QT, Trader} \iqfirm{toutes} \iqlevel{junior}]\label{iq:iv1:random-walks-and-martingales:10}
À chaque pas, une marche se déplace de $+1$ avec la probabilité $0.5$, de $-1$ avec la probabilité $0.3$, et reste sur place avec la probabilité $0.2$. Partant de 0, quelle est la probabilité qu'elle atteigne $+5$ avant $-5$~? Combien de temps cela prend-il en moyenne~?
\end{question}

\begin{question}[$\star\star$ \iqroles{QT, Trader, Risque, QR} \iqfirm{teneur de marché, hedge fund, crypto} \iqlevel{junior}]\label{iq:iv1:random-walks-and-martingales:11}
Un cours monte d'un tick avec la probabilité $0.55$ ou descend d'un tick avec la probabilité $0.45$, de façon indépendante. Vous achetez, avec un objectif de profit à 20 ticks à la hausse et un stop à 10 ticks à la baisse. Quelle est la probabilité que l'objectif soit atteint en premier, et combien de pas la transaction prend-elle en moyenne~? Quel est votre gain espéré en ticks~? Avec une marche équitable, quel serait-il alors~?
\end{question}

\begin{question}[$\star\star$ \iqroles{QT, Trader, Risque, Analyste} \iqfirm{hedge fund, teneur de marché, banque} \iqlevel{junior}]\label{iq:iv1:random-walks-and-martingales:12}
Vous disposez de 100~\$ et d'un pari à parité que vous gagnez avec la probabilité $0.55$, répété autant de fois que vous le souhaitez. Vous misez toujours le même montant en dollars. Quelle est la probabilité que vous soyez un jour ruiné si vous misez 1~\$ à chaque fois~? 10~\$~? 50~\$~?
\end{question}

\begin{question}[$\star\star$ \iqroles{QR, QT, Trader} \iqfirm{toutes} \iqlevel{junior}]\label{iq:iv1:random-walks-and-martingales:13}
Une marche équitable part de 0. Quelle est la probabilité qu'elle reste strictement au-dessus de 0 pendant les 20 premiers pas~? Quelle est la probabilité qu'elle revienne pour la première fois en 0 exactement au pas 2~? Au pas 4~? Au pas 6~?
\end{question}

\begin{question}[$\star\star$ \iqroles{QR, QT, Trader, Analyste} \iqfirm{hedge fund, teneur de marché} \iqlevel{junior}]\label{iq:iv1:random-walks-and-martingales:14}
Vous estimez qu'un événement a une probabilité $1/3$. Chaque jour apporte une nouvelle information qui double ou divise par deux la cote de cet événement, et votre croyance est toujours la probabilité conditionnelle correcte. Quelle est la probabilité que votre croyance atteigne un jour $8/9$~? Pourquoi n'avez-vous besoin ni de la fiabilité des nouvelles ni du nombre de jours~?
\end{question}

\begin{question}[$\star\star$ \iqroles{QR, QT, Trader} \iqfirm{hedge fund, teneur de marché} \iqlevel{junior}]\label{iq:iv1:random-walks-and-martingales:15}
Deux jetons se trouvent sur un anneau de 12 positions, distants de 4 positions l'un de l'autre. À chaque pas, les deux jetons se déplacent en même temps, chacun indépendamment d'une case dans le sens horaire ou dans le sens antihoraire avec la même probabilité. En moyenne, combien de pas faut-il pour qu'ils se retrouvent sur la même position~?
\end{question}

\begin{question}[$\star\star\star$ \iqroles{QR, QT, Risque} \iqfirm{hedge fund, teneur de marché} \iqlevel{junior, confirmé}]\label{iq:iv1:random-walks-and-martingales:16}
Une marche équitable se déroule sur 20 pas. À quel instant sa dernière visite en 0 est-elle la plus probable, et quelles sont les probabilités qu'elle se trouve à l'instant 0, à l'instant 10 et à l'instant 20~? À la limite de long terme, quelle est la probabilité que la dernière visite tombe dans le premier quart de la durée~?
\end{question}

\begin{question}[$\star\star\star$ \iqroles{QR, QT, Trader} \iqfirm{hedge fund, teneur de marché} \iqlevel{junior, confirmé}]\label{iq:iv1:random-walks-and-martingales:17}
Trois joueurs possèdent $a$, $b$ et $c$ pièces. À chaque tour, un couple ordonné de joueurs est tiré uniformément au hasard parmi les six couples ordonnés, et le premier donne une pièce au second. En moyenne, combien de tours faut-il pour qu'un joueur n'ait plus aucune pièce, en partant de 10, 20 et 30 pièces~? Si la partie se poursuit entre les deux joueurs restants, avec la même règle, combien de temps faut-il en moyenne pour qu'un joueur détienne les 60 pièces~?
\end{question}

\begin{question}[$\star\star\star$ \iqroles{QR, QT, Risque, DS} \iqfirm{hedge fund, teneur de marché} \iqlevel{junior, confirmé}]\label{iq:iv1:random-walks-and-martingales:18}
Une urne contient au départ une boule rouge et une boule bleue. Chaque fois, vous tirez une boule au hasard, puis vous la remettez accompagnée d'une boule supplémentaire de la même couleur. Montrez que la proportion de boules rouges est une martingale. Donnez la meilleure majoration possible, sans simulation, de la probabilité que cette proportion atteigne un jour $3/4$ ou plus, ainsi qu'une minoration issue du comportement à long terme de la proportion. Quelle est, approximativement, la vraie valeur~?
\end{question}

\begin{question}[$\star\star\star$ \iqroles{QR, QT, Risque} \iqfirm{hedge fund, banque} \iqlevel{junior, confirmé}]\label{iq:iv1:random-walks-and-martingales:19}
Une marche équitable se déroule sur 100 pas. Quelle est la moyenne de son point le plus haut, en comptant le départ en 0~? Utilisez le principe de réflexion pour le ramener à la loi de la valeur finale, et comparez le résultat à une approximation continue.
\end{question}

\begin{question}[$\star\star\star$ \iqroles{QR, QT, Risque} \iqfirm{hedge fund, teneur de marché} \iqlevel{junior, confirmé}]\label{iq:iv1:random-walks-and-martingales:20}
Une marche se déplace de $+1$ avec la probabilité $0.6$ et de $-1$ avec la probabilité $0.4$. Le premier passage en $+10$ dure en moyenne 50 pas. Quel est l'écart type de ce temps~? Trouvez-le à partir d'une martingale, et non d'une récurrence.
\end{question}

\begin{question}[$\star\star\star$ \iqroles{QR, QT, QD} \iqfirm{toutes} \iqlevel{junior, confirmé}]\label{iq:iv1:random-walks-and-martingales:21}
Une marche équitable se déplace sur le réseau des entiers, en choisissant l'un des points voisins avec la même probabilité à chaque pas, en dimensions un, deux et trois. Dans quelles dimensions revient-elle à son point de départ presque sûrement~? Donnez l'argument, ainsi que la probabilité de retour en dimension trois, à deux décimales près.
\end{question}

\begin{question}[$\star\star\star$ \iqroles{QR, QT, Trader} \iqfirm{hedge fund, teneur de marché} \iqlevel{junior, confirmé}]\label{iq:iv1:random-walks-and-martingales:22}
Une marche monte de 2 avec la probabilité $1/3$ et descend de 1 avec la probabilité $2/3$. Partant de 0, vous vous arrêtez dès que vous atteignez 10 ou plus, ou $-5$ ou moins. Quelle est la probabilité de vous arrêter par le haut, et combien de pas faut-il en moyenne~? Pourquoi le tiers, réponse que vous obtiendriez pour des pas de 1, n'est-il pas tout à fait juste~?
\end{question}

\begin{question}[$\star\star\star\star$ \iqroles{QR, QT} \iqfirm{hedge fund, teneur de marché} \iqlevel{confirmé}]\label{iq:iv1:random-walks-and-martingales:23}
Une marche équitable part de 0 et s'arrête dès qu'elle atteint pour la première fois $+n$ ou $-n$. Sa durée moyenne vaut $n^2$. Quelle est sa variance~? Donnez une formule et la valeur pour $n = 10$.
\end{question}

\begin{question}[$\star\star\star\star$ \iqroles{QR, QT} \iqfirm{hedge fund, teneur de marché} \iqlevel{confirmé}]\label{iq:iv1:random-walks-and-martingales:24}
Une marche équitable part de 0 et poursuit sa course jusqu'à son premier retour en 0. Elle revient presque sûrement, bien que le temps moyen soit infini. En moyenne, combien de fois visite-t-elle le point $+k$ avant de revenir, pour $k = 1$~? Pour $k = 100$~? Pourquoi cela n'est-il pas absurde~?
\end{question}

\begin{question}[$\star\star\star\star$ \iqroles{QR, QT, Trader, Risque} \iqfirm{hedge fund, teneur de marché} \iqlevel{confirmé}]\label{iq:iv1:random-walks-and-martingales:25}
Vous disposez de 25~\$ et vous avez besoin de 100~\$. Vous pouvez miser n'importe quel nombre entier de dollars à chaque tour d'un pari à parité que vous gagnez avec la probabilité $0.4$, et vous vous arrêtez dès que vous atteignez 100~\$ ou 0~\$. Quelle est votre probabilité de réussite si vous misez 1~\$ à la fois~? Et si vous misez à chaque fois le maximum possible sans dépasser 100~\$~? Est-ce la meilleure stratégie~? Qu'est-ce qui change si la probabilité de gain vaut $0.6$ ou $0.5$~?
\end{question}
